\subsection{The displayed restriction map does not lift categorically}
\label{sec:cyclic-eight-lift-obstruction}
The distinction between a positive ring map and a central tensor functor
is essential in Proposition~\ref{prop:cyclic-eight-center}.

\begin{proposition}[Obstruction for the specified cyclic-eight data]
\label{prop:cyclic-eight-lift-obstruction}
Let $\mathcal D$ be the actual cyclic-eight quotient category used in
Proposition~\ref{prop:cyclic-eight-center}, and let $f$ be its displayed
positive based-ring map. There is no fusion category $\mathcal C$ with
based ring $\RR{3}{215}$ and a braided equivalence
$Z(\mathcal D)\simeq Z(\mathcal C)$ under which the central forgetful
functor induces $f$. The same holds after dimension- and twist-preserving
permutations of the displayed central labels.
\end{proposition}
\begin{proof}
Write $d=4+\sqrt{17}$, so $d^2=8d+1$. The simple objects of
$\mathcal D$ are $\alpha^j,\alpha^j\rho$ for $j\in\ZZ/4\ZZ$, with
\[
 \rho\alpha=\alpha^{-1}\rho,
 \qquad \rho^2=1+2\sum_j\alpha^j\rho.
\]
The category has connected separable algebras $1+\alpha^j\rho$ for
every $j$; these come from the generalized Haagerup construction and
descend under the quotient tensor functor
\cite[Section~2.4 and Example~4.4]{GIcenters}.

Suppose the asserted equivalence and central functor exist, and let
$\mathcal M$ be the corresponding $\mathcal D$--$\mathcal C$ Morita
bimodule. The actual source restriction matrix $F_{\mathcal D}$ is
recorded in
\filename{research/center_continuation/cyclic8_canonical_restriction.json}.
Its complete tensor compatibility is verified against the same center
tensor as $f$. The trace identity given by the cross Gram matrix
$F_{\mathcal D}^{t}F$ yields
\[
 |\Irr(\mathcal M)|=6,\quad
 \operatorname{Tr}(L_\alpha)=0,\quad
 \operatorname{Tr}(L_{\alpha^2})=2,\quad
 \operatorname{Tr}(L_{\alpha^j\rho})=8,\quad
 \operatorname{Tr}(R_g)=2.
\]
Here $g$ is the nontrivial invertible of $\mathcal C$. The two columns
of $F$ used in these identities are constant on each class of central
simples having equal dimension and twist, so the permitted relabelings
do not change the argument.

The permutation $L_\alpha$ has one four-cycle and one two-cycle.
Symmetry and $L_\rho L_\alpha=L_\alpha^{-1}L_\rho$ put $L_\rho$ in
a reverse-circulant four-by-four block, an alternating cross block, and
a lower block $\left(\begin{smallmatrix}u&v\\v&u\end{smallmatrix}\right)$.
The trace equations give
$u=4-r_0-r_2$, $v=4-r_1-r_3$ and
$\operatorname{Tr}(L_\rho^2)=70$.
The exact enumeration in
\filename{research/new_frontier/cyclic8_rank6_modules.py}
checks all $2025$ bounded integer assignments. Exactly eight remain,
four of each type in the table below. Both cross entries equal one.
The squared dimensions use Morita normalization, so their sum is
$\operatorname{FPdim}(\mathcal D)=8+32d$.
\begin{center}
\begin{tabular}{@{}ccc@{}}
\toprule
Upper row sum & Four-cycle objects & Two-cycle objects\\
\midrule
$1$ & $1+d$ & $2+14d$\\
$7$ & $1+7d$ & $2+2d$\\
\bottomrule
\end{tabular}
\end{center}
The unique involutory permutation of trace two commuting with
$L_\alpha$ shifts its four-cycle by two and fixes its two-cycle.
This determines $R_g$.

The algebra $A=1+g$ exists because $g$ fixes a simple object. Indeed that
object gives a one-object module over the pointed subcategory generated
by $g$, trivializing its associator. The exact full Morita calculation
\filename{research/new_frontier/verify_frontier_full_morita.py}
shows that the based ring of
$\mathcal C'=A\text{-bimod}(\mathcal C)$ is again $\RR{3}{215}$.
It does not identify $\mathcal C$ and $\mathcal C'$ categorically.
In the second row of the table, pass to right $A$-modules in
$\mathcal M$. Each fixed simple has two module structures, whose
squared Morita dimensions are $(2+2d)/2=1+d$; each free pair gives one
simple of squared dimension $4(1+7d)/2=2+14d$.
Thus in either case there is a Morita bimodule, with right based ring
$\RR{3}{215}$, containing a simple $\iota$ with
$\iota\iota^*=1+\alpha^h\rho$. Rename $\alpha^h\rho$ as $\rho$.

For clarity, this reduction does not assume that every algebra
$1+\rho$ has the same module action. Its four $\alpha$-translates are
distinct. The upper $\rho$ block and its diagonal square equation force
both cross entries to be one. The remaining lower block has
$u+v=7$ and $|u-v|=1$. If $u=4$, the object $\iota^*\kappa$ in the
right category would have dimension $4d$ and endomorphism dimension
$5$. It has no invertible summands since $\iota$ and $\kappa$ have
different dimensions. Writing its other multiplicities as $a,b,c,e$
gives $e=c$, $a+b+2c=4$ and squared norm $a^2+b^2+2c^2$.
The possible norms are $4,6,8,10,16$, never $5$. Hence $u=3$, $v=4$.

Choose a Morita object $\lambda$ in a second module category over
$\mathcal D$ with $\lambda\lambda^*=1+\alpha\rho$.
In the composed bimodule, Frobenius reciprocity gives two distinct
simples $\beta_0=\lambda^*\iota$ and
$\beta_2=\lambda^*\alpha^2\iota$, each of dimension $d+1$.
The object $\Gamma=\lambda^*\kappa$ has dimension $4d$, endomorphism
dimension $5$, and contains both $\beta$'s once. Its three other simple
summands have total dimension $2(d-1)$. One of them, $\gamma$, therefore
has dimension at most $2(d-1)/3$.

All dimensions in this composed bimodule lie in $K=\QQ(\sqrt{17})$:
the simple Perron eigenspace of the sum of the right action matrices
is defined over $K$, and the known dimension $d+1$ fixes its scale.
Thus the connected internal end $B=\gamma^*\gamma$ has dimension a
square in $K$, at most $L=4(d-1)^2/9$. It cannot be the unit algebra;
otherwise the bimodule would be the regular module and its two distinct
simples of dimension $d+1$ would contradict the dimensions of
$\RR{3}{215}$.

The unit has multiplicity one in $B$ and the other invertible has
multiplicity at most one. The complete list of dimensions at most $L$ is
\[
 1,2,d,d+1,d+2,d+3,
 2d-1,2d,2d+1,2d+2,2d+3,2d+4,3d-2.
\]
Here $2$ is not a square in $K$, and the field norms of the other
nontrivial entries are respectively
\[
 -1,8,19,32,-19,-4,13,32,53,76,-53,
\]
none a nonnegative rational square. This contradicts the existence of
$B$. The exhaustive bound and these norms are verified by
\filename{research/new_frontier/cyclic8_small_algebra_dimensions.py}.
This extends the second-algebra argument of
\cite[Lemmas~4.15--4.16]{GISAH}; the additional possibility $3d-2$,
omitted from that paper's small-dimension list, has norm $-53$ and
does not affect the contradiction.
\end{proof}

The source category and specified induction data in this proposition
are substantive hypotheses. An abstract modular datum or an integral
Verlinde tensor does not identify a braided category with $Z(\mathcal D)$.
Consequently this proposition alone does not resolve general
categorifiability. The independent stabilizer-rotation argument in
Theorem~\ref{thm:stabilizer-rotation-frontier} supplies that general
exclusion without identifying an unknown center with $Z(\mathcal D)$.
